\documentclass[12pt]{article}

\usepackage[utf8]{inputenc}
\usepackage[margin=1in]{geometry}
\usepackage{setspace}
\usepackage{url}
\usepackage{mathptmx}

\doublespacing
\setlength{\emergencystretch}{3em}
\setlength{\parindent}{0pt}
\setlength{\parskip}{0.25em}

\begin{document}

\begin{center}
\textbf{Proposal for Comparative Essay}
\end{center}

Ma Toan Bach\\
Student Number: 112708227\\
LSO 100--Canadian Short Story (KCC)\\
Date: July 17, 2026

\vspace{0.5em}
\textbf{Selected Short Stories}

I plan to compare Alice Munro's ``Boys and Girls'' and Madeleine Thien's ``Simple Recipes.'' I chose these stories because both use first-person narration to explore how formative family experiences shape identity. Although the stories have different cultural settings, both examine the conflict between family love, obedience, and personal identity.

\textbf{Tentative Title}

\textit{Family Pressure, Identity, and Belonging in ``Boys and Girls'' and ``Simple Recipes''}

\textbf{Central Focus / Preliminary Thesis}

Both ``Boys and Girls'' and ``Simple Recipes'' show that family expectations strongly influence how young people understand themselves; however, Munro emphasizes the limiting effects of traditional gender roles, while Thien focuses on cultural conflict, violence, and divided family loyalty. Together, the stories suggest that family belonging can provide love and identity while also creating pressure, resentment, and emotional pain.

\textbf{Main Points of Comparison}

My first point of comparison will be the influence of the father in each story. In ``Boys and Girls,'' the narrator admires her father and wants to participate in his outdoor work, but she gradually learns that he does not view her as equal to a boy. In ``Simple Recipes,'' the narrator's father first represents care, tradition, and cultural connection through preparing rice, but his violence later damages that sense of security.

My second point will be identity and belonging within the family. Munro's narrator resists the domestic role expected of her because she does not want her gender to determine her future. Thien's narrator experiences divided loyalty because she loves her father while also feeling shame and fear after witnessing his treatment of her brother. I will also compare their retrospective narration: Munro's narrator looks back on her growing awareness of gender inequality, while Thien's narrator revisits a traumatic memory to understand her conflicting love and fear toward her father.

\textbf{Preliminary Evidence}

For ``Boys and Girls,'' I plan to discuss the narrator's pride in helping her father, her decision to open the gate for Flora, and her father's statement that she is ``only a girl.'' These events reveal her resistance to gender expectations and her realization that her family judges her according to her gender. For ``Simple Recipes,'' I will examine the rice-making ritual, the brother's rejection of the meal, and the violence at the dinner table. These moments show the movement from family closeness to conflict and explain the narrator's complicated feelings toward her father.

\textbf{Preliminary Sources}

Nischik's analysis of gender in ``Boys and Girls'' will support my discussion of the narrator's resistance to traditional female roles. Ty's chapter on East and Southeast Asian Canadian literature will provide cultural context for examining identity, immigration, and belonging in Thien's work. Together, these sources will connect the stories to broader discussions of gender and cultural identity in Canadian literature.

\clearpage
\begin{center}
\textbf{References}
\end{center}
\raggedright

\hangindent=0.5in\noindent Munro, A. (1968). \textit{Dance of the happy shades}. Ryerson Press.

\hangindent=0.5in\noindent Nischik, R. M. (2007). (Un-)doing gender: Alice Munro, ``Boys and Girls'' (1964). In R. M. Nischik (Ed.), \textit{The Canadian short story: Interpretations} (pp. 203--218). Camden House. \url{https://doi.org/10.1515/9781571136886-016}

\hangindent=0.5in\noindent Thien, M. (2001). \textit{Simple recipes}. McClelland \& Stewart.

\hangindent=0.5in\noindent Ty, E. (2016). (East and Southeast) Asian Canadian literature: The strange and the familiar. In C. Sugars (Ed.), \textit{The Oxford handbook of Canadian literature} (pp. 564--582). Oxford University Press. \url{https://doi.org/10.1093/oxfordhb/9780199941865.013.31}

\end{document}
